% Generated draft: requires mathematical and source review.
\documentclass[11pt]{article}
\usepackage[margin=25mm]{geometry}
\usepackage{fontspec}
\setmainfont{Latin Modern Roman}
\usepackage{amsmath,amssymb,mathtools,unicode-math}
\setmathfont{Latin Modern Math}
\usepackage{enumitem,needspace,xurl,hyperref}
\setlength{\parindent}{0pt}
\setlength{\parskip}{5pt}
\setlength{\emergencystretch}{3em}
\newcommand{\N}{\mathbb{N}}
\newcommand{\Z}{\mathbb{Z}}
\newcommand{\Q}{\mathbb{Q}}
\newcommand{\R}{\mathbb{R}}
\newcommand{\C}{\mathbb{C}}
\newcommand{\F}{\mathbb{F}}
\newcommand{\A}{\mathbb{A}}
\newcommand{\PP}{\mathbb{P}}
\newcommand{\E}{\mathbb{E}}
\newcommand{\eps}{\varepsilon}
\newcommand{\dd}{\,\mathrm{d}}
\begin{document}
\section*{UnsolvedMath ID: 30002256}
\textbf{Source problem code:} catalogue entry
\subsection{Definitions and mathematical statement}
Let \(G\) be a group and \(H\leq G\) a subgroup. Write \(\operatorname{Aut}(H)\) for the automorphism group of \(H\), \(\operatorname{Inn}(H)\) for its subgroup of inner automorphisms, and
\[
\operatorname{Out}(H)=\operatorname{Aut}(H)/\operatorname{Inn}(H).
\]
Define
\[
\operatorname{Aut}(G;H)=\{\Phi\in\operatorname{Aut}(G):\Phi(H)=H\}.
\]
Restriction to \(H\), followed by passage to outer automorphism classes, gives a homomorphism
\[
\operatorname{Aut}(G;H)\longrightarrow\operatorname{Out}(H),
\qquad \Phi\longmapsto[\Phi|_H].
\]
The group of extendable outer automorphisms is its image:
\[
\operatorname{Out}(H\hookrightarrow G)
=\bigl\{[\alpha]\in\operatorname{Out}(H):
\exists\Phi\in\operatorname{Aut}(G),\ 
\Phi(H)=H,\ \Phi|_H=\alpha\bigr\}.
\]
Here extendability means that an automorphism representative extends to an automorphism of \(G\). This condition is independent of the representative: conjugation by an element of \(H\) extends to conjugation by the same element of \(G\).

\textbf{Problem.} For a subgroup inclusion \(H\leq G\), describe the elements of \(\operatorname{Out}(H\hookrightarrow G)\), and determine whether this group is finitely generated or finitely presented. Finite generation means that some finite subset generates the group; finite presentability means that the group admits a presentation with finitely many generators and finitely many defining relations.

\textbf{Unresolved scope.} The supplied statement imposes no finiteness, geometric, or other hypotheses on \(G\) or \(H\). It also does not specify what kind of description is intended by ``what does this group contain.'' Thus the statement records a general characterization and finiteness question, without asserting a universal affirmative answer or specifying a narrower class of inclusions.
\subsection{Short English statement}
For a subgroup \(H\leq G\), consider the outer automorphism classes of \(H\) represented by restrictions of automorphisms of \(G\) preserving \(H\). Describe this subgroup of \(\operatorname{Out}(H)\) and investigate its finite generation and finite presentability. The supplied question leaves the classes of groups and the intended form of the description unspecified.
\subsection{Sources}
\begin{itemize}
\item Supplied local UnsolvedMath registry record: permanent ID \(30002256\), source problem code \texttt{catalogue entry}, title \emph{Extendable Outer Automorphisms of Subgroup Inclusions}; registry locations \texttt{data/unsolvedmath/problems.json} and \texttt{catalog/problems.json}. The record assigns the year 2013 and supplies no individual proposer.
\item \emph{Model Theory: Groups, Geometry, and Combinatorics} (2013), Oberwolfach Reports, pp.\ 5--66, DOI \texttt{10.4171/owr/2013/01}: \url{https://doi.org/10.4171/owr/2013/01}. Also listed in the supplied literature sources as \emph{Model Theory, Oberwolfach Report 01/2013}. The supplied background describes Oberwolfach Reports as documenting mathematical workshops and their open-problem discussions.
\item \emph{Splittings and Automorphisms of Relatively Hyperbolic Groups}: \url{https://arxiv.org/abs/1212.1434}. The supplied assessment attributes to Guirardel and Levitt a description of the extendable subgroup and finite-index McCool and finite-presentation results for malnormal quasiconvex and relatively hyperbolic cases, but not for arbitrary subgroup inclusions. The precise hypotheses and conclusions of these results are not included in the record.
\item Literature assessment date: \texttt{2026-08-21}. The supplied dated curation labels the problem \texttt{partially\_solved}. This assessment, including the attributed literature claims, is not independently verified here; the source itself requests independent verification of time-sensitive status and cited claims.
\item The supplied statement status is \texttt{corrected\_verified}, with a supplied verification note saying that the official report states the general question with the same extendability definition and finiteness alternatives. This is registry-provided source verification, not an independent verification performed for this draft.
\end{itemize}
\end{document}
